\documentclass[11pt,a4paper]{article}
\usepackage[margin=1in]{geometry}
\usepackage{enumitem}
\usepackage{titlesec}
\usepackage{xcolor}
\usepackage{hyperref}
\usepackage{parskip}

\titleformat{\section}{\large\bfseries}{\thesection}{1em}{}
\titleformat{\subsection}{\normalsize\bfseries}{\thesubsection}{1em}{}

\title{\textbf{GramMitra\_AI \\ \large Multilingual AI Farmer Assistant \\ \large Project Scope \& Team Collaboration Reference}}
\author{Final Year Computer Engineering Project \\ ISBM College of Engineering, Pune (SPPU)}
\date{}

\begin{document}
\maketitle

\section{Problem Statement}
Farmers in India face two connected gaps. First, agricultural advisory --
crop selection, fertilizer use, weather-based decisions, and awareness of
government schemes -- is often not available in their own language.
Second, when something goes wrong (electricity outage affecting
irrigation, water supply issues, crop damage), there is no simple way to
know which government department to approach. This project builds a
single multilingual AI-assisted platform that gives farmers proactive
agricultural guidance in their own language, connects them to relevant
government scheme information, and helps them correctly route complaints
to the right department.

\section{Scope}

\subsection{In Scope (Version 1)}
\begin{itemize}[leftmargin=*]
  \item \textbf{Crop recommendation} -- soil type, N-P-K values, and
    rainfall/region as input, suggested crops as output (ML model trained
    on a crop-recommendation dataset).
  \item \textbf{Weather-based advisory} -- live weather data via API,
    translated into simple actionable advice for farming decisions.
  \item \textbf{Fertilizer suggestion} -- crop and soil deficiency mapped
    to a fertilizer type/quantity recommendation.
  \item \textbf{Multilingual AI chatbot} -- farmer queries answered in
    English and at least one regional language (Hindi/Marathi), using a
    pretrained multilingual NLP model (e.g. IndicBERT) rather than a
    model trained from scratch.
  \item \textbf{Complaint classification \& routing} -- farmer describes
    an issue (e.g. electricity, water, crop damage); the system classifies
    the category and returns the correct department, contact information,
    and a reference ID. \emph{Note: this is a smart routing
    recommendation within our own system, not a live integration with
    government complaint-filing APIs, which are not publicly available.}
  \item \textbf{Government scheme \& subsidy guidance} -- folded into the
    chatbot's knowledge base; answers eligibility and application
    questions on schemes such as PM-KISAN, PMFBY (crop insurance),
    PM-KUSUM, and Soil Health Card, in the farmer's own language.
  \item \textbf{Live market price checker} -- mandi/crop prices pulled
    from a public source (e.g. Agmarknet/e-NAM open data via
    data.gov.in), filterable by crop and district.
  \item \textbf{Progressive Web App (PWA)} -- installable, offline-caching
    layer added once the core web platform is stable, to support farmers
    on low-end phones with patchy connectivity.
\end{itemize}

\subsection{Out of Scope}
\begin{itemize}[leftmargin=*]
  \item Real integration with government complaint-filing or insurance
    claim-status APIs (not publicly available to student developers).
  \item Native mobile app (PWA covers the mobile use case; native app is
    only a stretch goal if there is real time remaining).
  \item Payment or subsidy disbursement processing.
  \item IoT/hardware sensors -- all inputs are entered manually by the
    farmer.
  \item Cold-storage/warehouse locator and credit/loan processing
    (considered and explicitly excluded -- no reliable open dataset for
    the former, regulatory territory for the latter).
\end{itemize}

\subsection{Stretch Goals (only after core features are stable)}
\begin{itemize}[leftmargin=*]
  \item Voice input for the chatbot.
  \item Farmer login with query/complaint history dashboard.
  \item SMS/notification alerts for weather warnings.
  \item Native mobile app built on top of the PWA.
\end{itemize}

\subsection{Scope Change Policy}
\begin{itemize}[leftmargin=*]
  \item \textbf{Feature-level changes} (adding/removing/swapping a
    feature) require agreement from all team members, not a single
    developer's decision. Propose changes as a GitHub Issue tagged
    \texttt{scope-change} with a one-line reason \emph{before} any code
    is written, and update this document in the same pull request once
    agreed.
  \item \textbf{Implementation-level decisions} (e.g. choice of model,
    library, internal code structure) are the responsibility of whoever
    owns that module -- no team vote required.
  \item \textbf{Scope freeze:} the feature list locks around month 4 of
    the project timeline. After the freeze, only bug fixes and polish are
    permitted -- no new features or swaps.
\end{itemize}

\section{Team Collaboration \& Git Workflow}

\subsection{Repository Structure}
\begin{verbatim}
frontend/   React web app
backend/    API server (FastAPI/Flask)
ml/         ML models: crop recommendation, fertilizer suggestion,
            complaint classifier, scheme/price lookup logic
docs/       Architecture notes, API contract, this scope document
\end{verbatim}

\subsection{Branching Strategy}
\begin{itemize}[leftmargin=*]
  \item \texttt{main} -- always working, demo-ready. Protected; no direct
    pushes; changes land only via a reviewed pull request.
  \item \texttt{develop} -- integration branch; all feature branches
    merge here first.
  \item \texttt{feature/<short-name>} -- one branch per task, e.g.
    \texttt{feature/crop-recommendation}, branched from \texttt{develop}.
  \item \texttt{fix/<short-name>} -- bug fixes.
\end{itemize}

\subsection{Divide the Work -- Module Ownership}
With four team members, each person owns one clear module boundary so
work does not overlap on the same files:
\begin{itemize}[leftmargin=*]
  \item \textbf{Frontend} -- UI, multilingual interface, PWA setup.
  \item \textbf{Backend/API} -- endpoints, database, complaint routing
    and scheme/price lookup logic.
  \item \textbf{ML/NLP} -- crop recommendation model, fertilizer logic,
    chatbot NLP integration, complaint classifier.
  \item \textbf{Data, integrations \& documentation} -- datasets, weather
    and market-price API integration, keeping this document and the
    report in sync with what is actually built.
\end{itemize}
Roles can overlap slightly so no single person is a bottleneck, but each
module should have one clear owner.

\subsection{Contract-First Development}
Before writing code for any feature, define the exact API request/response
shapes and database schema in \texttt{docs/api-contract.md} as a team.
This is what allows four people (or their AI coding tools) to build
independently and still have the pieces fit together -- it does not
matter that different developers or different AI tools produce different
code, as long as everyone builds against the same agreed contract.

\subsection{Shared Conventions}
\begin{itemize}[leftmargin=*]
  \item One shared starter repository skeleton -- everyone builds on top
    of it rather than generating a full app independently from scratch.
  \item Shared linting/formatting (ESLint + Prettier for JS, Black +
    Flake8 for Python) enforced via a pre-commit hook, so code is
    normalized regardless of how it was generated.
  \item A short conventions document (naming, folder structure, error
    handling) pasted as context whenever prompting an AI coding tool, so
    different people's AI-assisted output converges on the same
    patterns.
  \item Commit messages: short and specific (e.g. ``Add crop
    recommendation API endpoint''), not vague (``update'', ``fix
    stuff'').
\end{itemize}

\subsection{Code Review \& Integration}
\begin{itemize}[leftmargin=*]
  \item Every pull request into \texttt{develop} requires at least one
    teammate's review before merging -- this catches logic that does not
    match the agreed contract, not just style issues.
  \item Basic integration tests check that each API endpoint returns the
    fields the frontend expects, catching contract drift early.
  \item A fixed weekly integration sync -- merge everyone's branches into
    \texttt{develop} and confirm the whole app still runs together,
    rather than discovering integration issues in the final weeks.
\end{itemize}

\subsection{Secrets \& Environment}
Never commit API keys, database passwords, or \texttt{.env} files -- use
a \texttt{.gitignore} from the start and share secrets separately
(e.g. a shared password manager), never inside the repository.

\section{Suggested Timeline (6--8 Months)}
\begin{itemize}[leftmargin=*]
  \item \textbf{Months 1--2:} Literature review, dataset finalization,
    architecture design, backend + crop/weather/fertilizer core built.
  \item \textbf{Months 3--4:} Multilingual chatbot, complaint routing,
    scheme guidance, and market price checker built; git workflow and
    weekly integration running. Scope freeze at end of month 4.
  \item \textbf{Month 5:} Stabilize the web platform end-to-end; begin
    PWA conversion once the frontend is stable.
  \item \textbf{Month 6:} PWA polish, testing, evaluation metrics
    gathered.
  \item \textbf{Months 7--8:} Buffer for bug fixes, documentation, report
    writing, and presentation preparation.
\end{itemize}

\end{document}
